% Canonical simulation protocol for the linguistic-evolution toolkit.
% Protocol version: corrected v2. Last verified against the repository on
% 2026-08-19. This file intentionally contains no document preamble so that it
% can be uploaded to Overleaf and included with:
%     % Canonical simulation protocol for the linguistic-evolution toolkit.
% Protocol version: corrected v2. Last verified against the repository on
% 2026-08-19. This file intentionally contains no document preamble so that it
% can be uploaded to Overleaf and included with:
%     % Canonical simulation protocol for the linguistic-evolution toolkit.
% Protocol version: corrected v2. Last verified against the repository on
% 2026-08-19. This file intentionally contains no document preamble so that it
% can be uploaded to Overleaf and included with:
%     % Canonical simulation protocol for the linguistic-evolution toolkit.
% Protocol version: corrected v2. Last verified against the repository on
% 2026-08-19. This file intentionally contains no document preamble so that it
% can be uploaded to Overleaf and included with:
%     \input{canonical_simulation_protocol}
% Change \section to \subsection (etc.) if embedding below another heading.
% Requires: amsmath, natbib, booktabs, and array (all already loaded by the
% current Overleaf manuscript). The bibliography should contain berg1995.

\section{Canonical simulation protocol (corrected v2)}
\label{sec:canonical-simulation-protocol}

\paragraph{Scope and precedence.}
This section is the canonical description of the corrected-v2 simulation
workflow. It covers the current memory-primary, informed-noise experiments,
including the completed task-order comparison and the configured population,
identity, and scripted-defector follow-ups. Earlier project phases used
materially different noise and memory implementations and should not be
described using this protocol. In particular, the current implementation keeps
separate actual and communicated ledgers: communication noise changes what an
agent observes about another agent's action, but it does not change the actual
transfer or the ground-truth payoff calculation.

\subsection{System overview}
\label{sec:protocol-overview}

Each simulation run consists of a population of same-model LLM agents playing
ten rounds of a repeated trust game, optionally interleaved with a myth-writing
task. The simulation is centrally orchestrated. Agents do not exchange
free-form chat messages. Their inter-agent information channels are (i) game
actions and communicated outcomes, (ii) any explicitly supplied game-history
or reputation block, and (iii) myths transmitted between rounds. Each agent
has a private rolling chat history; no agent receives another agent's private
conversation memory.

An experiment set is defined in a YAML configuration and expanded into the
Cartesian product of model, system-prompt template, persona, task order, myth
topic, game-parameter preset, myth-prompt arm, and replicate. Each expanded
combination becomes one independent simulation run. Replicates may be executed
in separate worker processes, while the state transitions within any one run
remain ordered. In seeded follow-up protocols, the exogenous pairing and noise
processes are reproducible conditional on their recorded seeds; model sampling
remains stochastic.

The three primary task orders are:
\begin{itemize}
    \item \textbf{Game only}: one game interaction per agent per round and no
    myth-writing task.
    \item \textbf{Game $\rightarrow$ myth}: agents make the current game
    decision before writing the current-round myth. The new myth can therefore
    affect later rounds, but not the game decision that preceded it.
    \item \textbf{Myth $\rightarrow$ game}: agents write the current-round myth
    before making the current game decision. Because both tasks use the same
    private memory, the new myth is available to the immediately following game
    call.
\end{itemize}
Whenever a run contains both tasks, every current game prompt contains exactly
once the instruction that the agent should take myths written in the session
into account when deciding. This explicit link is held constant across the two
myth-present task orders.

\subsection{Agents, prompts, and private memory}
\label{sec:protocol-agents-memory}

All agents in a run share the configured base model, temperature, neutral
persona, and prompt templates, but each agent has a distinct identifier,
optional display name, system prompt, and private message buffer. Unless a
condition states otherwise, current experiments use Claude Sonnet 4.5,
temperature $0.8$, and a neutral persona. The protected confirmatory launcher
routes these calls directly to Anthropic and fixes the maximum response length
at 4,096 tokens.

The system prompt explains both roles of the trust game, the payoff rules, the
required decision format, and the possible myth-writing task. In informed-noise
conditions it also states that communicated amounts may differ slightly from
the amounts actually sent or returned. Population-setting text tells agents
whether partners are fixed or rotating and whether identities are visible.

Memory capacity is measured in completed user--assistant exchanges, not in raw
messages. The system message is always preserved. Before a new model call, the
buffer retains at most the configured number of preceding exchanges and then
adds the current user prompt. In the canonical \emph{memory-primary} regime,
accepted game and myth exchanges are written to the same private buffer. Game-
only cells use capacity three. Two-task cells use capacity six, so both regimes
retain approximately three completed rounds of task interaction: three game
exchanges in the former and roughly three game-plus-myth pairs in the latter.
This matching prevents the addition of the myth task from mechanically
shortening the intended round-level memory horizon.

The exact request and response for every interaction are also copied to a
non-truncated audit history. This audit record includes the prompt, the full
message list supplied to the provider, the response text, available reasoning
and token usage, task and round metadata, role, opponent, dyad, response source,
and any error. The rolling private memory controls agent behaviour; the audit
history is retained only for provenance and verification.

\subsection{Population structure and pairings}
\label{sec:protocol-pairings}

Every round partitions the population into sender--receiver dyads. Within a
dyad, the sender always acts before the receiver. Pairing metadata is saved at
both the round and dyad levels.

For two-agent runs, the same pair remains together for all rounds and sender
and receiver roles alternate by round. For populations larger than two, two
pairing modes are available:
\begin{itemize}
    \item \textbf{Balanced rotating}: a full-run schedule balances sender and
    receiver assignments while preferentially pairing agents who have met less
    often. A recorded pairing seed makes the schedule reproducible.
    \item \textbf{Fixed}: adjacent agents form persistent dyads and roles
    alternate within each dyad. Thus an eight-agent fixed run consists of four
    parallel repeated dyads.
\end{itemize}
Names may be shown or hidden. When names are hidden, prompts use relational
labels such as ``your current co-player'' rather than stable identity labels.
The \texttt{self\_and\_coplayer} history policy can additionally insert a
bounded summary of earlier games involving the current co-player. This is
reputation information assembled from the public simulation record; it is not
access to that co-player's private LLM context.

The completed eight-agent task-order comparison uses balanced rotating
partners, hidden names, and up to three earlier games involving the current
co-player. The completed two-agent comparison is a repeated dyad. These arms
therefore differ simultaneously in population size, partner rotation,
anonymity, and supplied reputation information. They must be described as a
\emph{population-regime comparison}, not as a pure population-size
manipulation. The configured population-isolation follow-up removes these
differences by comparing one fixed dyad with four fixed dyads under unified
prompts, hidden names, and no injected history.

\subsection{Repeated trust game}
\label{sec:protocol-trust-game}

We use a role-swapping repeated trust game \citep{berg1995} with endowment
$E=\$5$ and multiplier $m=3$. In dyad $d$ and round $t$, the sender selects an
actual transfer $s_{td}\in[0,E]$. The receiver's actual allocation is
$m s_{td}$ and the receiver selects an actual return
$r_{td}\in[0,m s_{td}]$. The ground-truth per-round payoffs are
\begin{equation}
    \pi^{S}_{td}=E-s_{td}+r_{td},
    \qquad
    \pi^{R}_{td}=m s_{td}-r_{td}.
    \label{eq:trust-game-payoffs}
\end{equation}
Actual cumulative balances are the sums of these payoffs over the rounds in
which an agent participates.

The sender is prompted first and returns a JSON-like object containing a
\texttt{send} amount. The parsed amount is bounded to $[0,E]$ and stored before
the receiver prompt is constructed. The receiver is then shown the current
communicated send, responds with a \texttt{return} amount, and that amount is
bounded to the feasible actual interval $[0,m s_{td}]$. The complete response,
the parsed action, any clamping, payoffs, and cumulative balances are retained.
A decision response that cannot be parsed causes the run to fail closed rather
than silently inventing an action.

\subsection{Communication noise and the two ledgers}
\label{sec:protocol-noise}

In the primary corrected-v2 condition, both sends and returns are transmitted
through informed bidirectional uniform communication noise with range $a=\$1$.
For an action $x$ with feasible upper bound $u$, the communicated amount is
\begin{equation}
    \widetilde{x}
    = \operatorname{clip}\!\left(x+\epsilon,0,u\right),
    \qquad
    \epsilon\sim\mathcal{U}(-a,a),
    \label{eq:communication-noise}
\end{equation}
rounded to two decimal places. Noise is generated separately for the send and
return events. With seeded protocols, the random draw is a deterministic
function of the noise seed, round, dyad identifier, and action type. This makes
the exogenous communication channel reproducible without coupling it to model
sampling.

The environment maintains parallel actual and communicated views:
\begin{itemize}
    \item The receiver observes $\widetilde{s}_{td}$ and is told that its
    visible receipt is $m\widetilde{s}_{td}$, while the ground-truth allocation
    remains $m s_{td}$.
    \item The sender later observes $\widetilde{r}_{td}$ rather than
    $r_{td}$.
    \item Actual payoffs and balances use Equation~\ref{eq:trust-game-payoffs}.
    The sender's visible payoff is $E-s_{td}+\widetilde{r}_{td}$, and the
    receiver's visible payoff is $m\widetilde{s}_{td}-r_{td}$. Parallel visible
    balances accumulate these quantities for use in later prompts.
\end{itemize}
Thus noise perturbs perceived intentions and feedback, not the underlying
economic transaction. Communicated sends are generated only after the current
actual send exists and while constructing the receiver's prompt. Missing
current transfers raise an error; they never fall back to zero.

\subsection{Myth writing and cultural transmission}
\label{sec:protocol-myths}

In myth-present conditions, every active agent writes one myth per round. In
round one, the prompt requests an approximately 200-word myth and allows any
mythic setting, characters, or symbols. It states that the myth and game are
connected and asks for a story reflecting how the game should be played. In
later rounds, the memory-primary prompt quotes the myth written by the agent's
co-player in the preceding round. It asks the agent to write a new myth that
adapts its earlier work and continues to reflect how the game should be played.
The agent's own earlier myths are available through its private message buffer
rather than duplicated in the current prompt.

Myth transmission therefore has a one-round lag. In a rotating population, an
agent receives the previous round's partner myth, which need not have been
written by its current game partner. The myth channel consequently propagates
text along the realised pairing network rather than creating a population-wide
broadcast. All agents' myth calls within a round are issued concurrently after
their prompts have been prepared, so no agent sees another agent's myth from
the same round.

Before a myth response is committed to memory or the simulation state, a
validator checks that it is non-empty and does not appear to continue into a
game prompt. To avoid false positives from ordinary story language, rejection
requires at least two independent prompt-continuation markers. A rejected myth
is retained in the audit trail, the unanswered prompt is removed from private
memory, and the call is retried once from the same clean conversation state. If
the retry is also invalid, the run stops and writes an error artifact.

\subsection{Scripted-defector and monitor extensions}
\label{sec:protocol-extensions}

The scripted-defector treatment assigns a reproducible subset of agents to a
\texttt{forced\_zero} game policy. A scripted defector always sends zero as a
sender and returns zero as a receiver. These game actions do not consume an LLM
call, but the corresponding prompt and scripted response may still be written
to the agent's private memory so its subsequent myth writing can respond to its
experience. Defectors use the normal LLM myth-writing path. In the planned
primary treatment, two of eight agents are defectors, agent names are hidden,
and the treatment role is not disclosed to the agents.

The engine also contains an optional silent myth monitor. When explicitly
enabled, a separate model flags myths that communicate actionable game
strategy, and a flagged agent's current-round game earnings are removed from
its cumulative balance without a textual notification. The underlying
transaction and payoff records remain unchanged. This monitor is not part of
the corrected-v2 task-order comparison and must not be assumed active unless a
run's metadata records an enabled monitor configuration.

\subsection{Execution, retries, and persistence}
\label{sec:protocol-execution}

At launch, the batch runner records the source commit, whether the worktree is
dirty, the absolute configuration path and SHA-256 digest, provider route,
resolved provider model, maximum output-token setting, task order, prompt arm,
population settings, memory policy, and all explicit seeds. The protected
confirmatory launcher refuses to start from a dirty worktree, uses a fresh
output directory, requires the expected number of completed runs in every
cell, and audits each cell before proceeding.

Provider adapters make up to three attempts for errors they classify as
retryable, with exponential backoff. For the direct Anthropic route used in the
confirmatory batch, this includes rate-limit, connection, timeout, server, and
empty-response failures. If all provider attempts fail, any unanswered user
prompt is removed from the private buffer and the failed attempt remains in the
audit history. Myth validation has the additional single clean retry described
above. Scripted game actions are recorded with a distinct response source so
they are never counted as LLM calls.

After every successful round, the runner atomically writes a lightweight
\texttt{.results.json} containing the simulation state without agent message
histories. At the configured checkpoint interval (ten rounds in the current
batch runners), it also writes a full ordinary checkpoint containing the agent
buffers. If that ordinary checkpoint exists, a subsequent launch resumes at
the first unfinished round. On an exception, the runner writes a separate full
\texttt{.error.json} artifact for diagnosis before re-raising the error; this
error artifact is not selected automatically as a resume point. On successful
completion the runner writes (i) a full-state \texttt{.json}, (ii) the
lightweight results file, (iii) a verbatim \texttt{.log}, and (iv) a
human-readable \texttt{.transcript.pdf}, then removes the ordinary checkpoint.
Output paths are organised by output batch, experiment set, model, task order,
and named game-parameter preset.

\subsection{Protocol audit}
\label{sec:protocol-audit}

Final outputs are admitted to substantive analysis only after a strict audit.
For every run, the audit verifies:
\begin{itemize}
    \item the configured number of rounds, agents, dyads, game responses, myths,
    and accepted interactions;
    \item the absence of unrecovered provider or validation errors and the
    correct accounting of recovered attempts;
    \item the informed-noise notice in system prompts and the numerical bounds
    on every communicated send and return;
    \item the presence of the myth-decision instruction exactly once in each
    applicable current game prompt and its absence from game-only prompts;
    \item private-memory lengths and the configured explicit-history windows;
    \item partner stability in fixed-pair conditions, population-setting text,
    name visibility, scripted-response counts, and defector disclosure policy;
    \item execution provenance, including the source revision, configuration
    digest, provider route, and output-token setting; and
    \item identical realised pairing schedules across runs that share a nominal
    paired-protocol key.
\end{itemize}
The completed corrected-v2 task-order batch contains 60 final runs (six cells,
ten runs per cell). All 60 passed the protocol audit. The completed cells did
not, however, record the intended pairing and noise seeds; their inferential
comparisons must therefore treat runs as independent rather than paired.

\subsection{Outcome measures and statistical unit}
\label{sec:protocol-outcomes}

The independent statistical unit is a complete simulation run, not an
individual dyadic transaction or model call. Primary behavioural summaries are
computed from the actual ledger:
\begin{align}
    \text{trust ratio} &= s_{td}/E, \\
    \text{return ratio} &=
    \begin{cases}
        r_{td}/(m s_{td}), & m s_{td}>0,\\
        0, & m s_{td}=0,
    \end{cases} \\
    \text{final balance} &= \text{mean final actual balance across agents}.
\end{align}
We additionally report mean sent and returned amounts, joint population
balance, and round-level trust and return trajectories. For the completed
six-cell confirmatory comparison, between-condition contrasts use independent-
group Welch tests and Welch confidence intervals, with Holm adjustment across
the six prespecified within-population balance contrasts. Population-regime
interactions are estimated as independent differences-in-differences. Seeded
follow-up designs should preserve replicate identifiers and the shared
exogenous pairing/noise schedules so that their paired structure can be used in
the corresponding analysis.

The linguistic outputs support complementary analyses of within-agent myth
retention, between-agent and population-level semantic convergence, lexical
and structural change, strategy content, and links between myth features and
subsequent game behaviour. Such analyses operate on the accepted myths saved in
the final state; rejected continuation attempts are audit artifacts and are
never treated as cultural outputs.

\subsection{Current experimental matrix and interpretation}
\label{sec:protocol-current-matrix}

Table~\ref{tab:current-protocol-matrix} separates completed evidence from
implemented follow-ups. This distinction is part of the protocol: configured
or smoke-tested cells must not be described as substantive results.

\begin{table}[t]
\centering
\small
\setlength{\tabcolsep}{3pt}
\caption{Current corrected-v2 experimental matrix as of 19 August 2026.}
\label{tab:current-protocol-matrix}
\begin{tabular}{p{0.22\linewidth}p{0.31\linewidth}p{0.15\linewidth}p{0.22\linewidth}}
\toprule
Comparison & Design & Status & Permitted interpretation \\
\midrule
Task order by population regime & Repeated two-agent dyad versus anonymous,
rotating eight-agent population with co-player reputation; game only,
game$\rightarrow$myth, and myth$\rightarrow$game & 60/60 complete and audited
& Task-order and population-regime effects; not a pure agent-count effect \\
\addlinespace
Population isolation & One fixed dyad versus four fixed dyads, unified prompts,
hidden names, and no injected history, crossed with all three task orders &
Implemented; no substantive batch & Cleaner population-size test with partner
structure held fixed \\
\addlinespace
Identity isolation & Eight-agent balanced rotating population, no injected
history, hidden versus stable visible names & Implemented; no substantive batch
& Effect of stable partner identity under rotation \\
\addlinespace
Scripted defectors & Eight-agent rotating population, 0\% versus 25\%
\texttt{forced\_zero} agents, crossed with all three task orders & Implemented
and smoke-tested; no substantive batch & Cultural and behavioural response to
exogenous defection \\
\bottomrule
\end{tabular}
\end{table}

\subsection{Repository sources of truth}
\label{sec:protocol-source-files}

The prose above is implemented by the following repository files:
\begin{itemize}
    \item experiment definitions and prompt text:
    \texttt{config/experiments\_noisy.yaml};
    \item configuration expansion, provenance capture, execution, and output
    layout: \texttt{experiments/run\_noisy\_batch.py};
    \item round orchestration, task ordering, memory modes, checkpoints, and
    state serialization: \texttt{src/simulation.py};
    \item private message buffers and interaction audit records:
    \texttt{src/agents.py};
    \item trust-game prompts, communication noise, actual/visible ledgers, and
    payoff calculations: \texttt{games/trust\_game\_noisy.py};
    \item population assignment, pairings, identity presentation, and scripted
    policies: \texttt{games/dyadic\_pairing.py};
    \item myth prompting, one-round transmission, and response validation:
    \texttt{src/myth\_writer.py};
    \item protected confirmatory launch:
    \texttt{scripts/run\_corrected\_v2\_confirmatory\_20260812.sh}; and
    \item admissibility audit: \texttt{scripts/audit\_v2\_protocol.py}.
\end{itemize}
If this prose and an older manuscript description disagree, the saved run
metadata and these implementation files determine what was actually executed.

% Change \section to \subsection (etc.) if embedding below another heading.
% Requires: amsmath, natbib, booktabs, and array (all already loaded by the
% current Overleaf manuscript). The bibliography should contain berg1995.

\section{Canonical simulation protocol (corrected v2)}
\label{sec:canonical-simulation-protocol}

\paragraph{Scope and precedence.}
This section is the canonical description of the corrected-v2 simulation
workflow. It covers the current memory-primary, informed-noise experiments,
including the completed task-order comparison and the configured population,
identity, and scripted-defector follow-ups. Earlier project phases used
materially different noise and memory implementations and should not be
described using this protocol. In particular, the current implementation keeps
separate actual and communicated ledgers: communication noise changes what an
agent observes about another agent's action, but it does not change the actual
transfer or the ground-truth payoff calculation.

\subsection{System overview}
\label{sec:protocol-overview}

Each simulation run consists of a population of same-model LLM agents playing
ten rounds of a repeated trust game, optionally interleaved with a myth-writing
task. The simulation is centrally orchestrated. Agents do not exchange
free-form chat messages. Their inter-agent information channels are (i) game
actions and communicated outcomes, (ii) any explicitly supplied game-history
or reputation block, and (iii) myths transmitted between rounds. Each agent
has a private rolling chat history; no agent receives another agent's private
conversation memory.

An experiment set is defined in a YAML configuration and expanded into the
Cartesian product of model, system-prompt template, persona, task order, myth
topic, game-parameter preset, myth-prompt arm, and replicate. Each expanded
combination becomes one independent simulation run. Replicates may be executed
in separate worker processes, while the state transitions within any one run
remain ordered. In seeded follow-up protocols, the exogenous pairing and noise
processes are reproducible conditional on their recorded seeds; model sampling
remains stochastic.

The three primary task orders are:
\begin{itemize}
    \item \textbf{Game only}: one game interaction per agent per round and no
    myth-writing task.
    \item \textbf{Game $\rightarrow$ myth}: agents make the current game
    decision before writing the current-round myth. The new myth can therefore
    affect later rounds, but not the game decision that preceded it.
    \item \textbf{Myth $\rightarrow$ game}: agents write the current-round myth
    before making the current game decision. Because both tasks use the same
    private memory, the new myth is available to the immediately following game
    call.
\end{itemize}
Whenever a run contains both tasks, every current game prompt contains exactly
once the instruction that the agent should take myths written in the session
into account when deciding. This explicit link is held constant across the two
myth-present task orders.

\subsection{Agents, prompts, and private memory}
\label{sec:protocol-agents-memory}

All agents in a run share the configured base model, temperature, neutral
persona, and prompt templates, but each agent has a distinct identifier,
optional display name, system prompt, and private message buffer. Unless a
condition states otherwise, current experiments use Claude Sonnet 4.5,
temperature $0.8$, and a neutral persona. The protected confirmatory launcher
routes these calls directly to Anthropic and fixes the maximum response length
at 4,096 tokens.

The system prompt explains both roles of the trust game, the payoff rules, the
required decision format, and the possible myth-writing task. In informed-noise
conditions it also states that communicated amounts may differ slightly from
the amounts actually sent or returned. Population-setting text tells agents
whether partners are fixed or rotating and whether identities are visible.

Memory capacity is measured in completed user--assistant exchanges, not in raw
messages. The system message is always preserved. Before a new model call, the
buffer retains at most the configured number of preceding exchanges and then
adds the current user prompt. In the canonical \emph{memory-primary} regime,
accepted game and myth exchanges are written to the same private buffer. Game-
only cells use capacity three. Two-task cells use capacity six, so both regimes
retain approximately three completed rounds of task interaction: three game
exchanges in the former and roughly three game-plus-myth pairs in the latter.
This matching prevents the addition of the myth task from mechanically
shortening the intended round-level memory horizon.

The exact request and response for every interaction are also copied to a
non-truncated audit history. This audit record includes the prompt, the full
message list supplied to the provider, the response text, available reasoning
and token usage, task and round metadata, role, opponent, dyad, response source,
and any error. The rolling private memory controls agent behaviour; the audit
history is retained only for provenance and verification.

\subsection{Population structure and pairings}
\label{sec:protocol-pairings}

Every round partitions the population into sender--receiver dyads. Within a
dyad, the sender always acts before the receiver. Pairing metadata is saved at
both the round and dyad levels.

For two-agent runs, the same pair remains together for all rounds and sender
and receiver roles alternate by round. For populations larger than two, two
pairing modes are available:
\begin{itemize}
    \item \textbf{Balanced rotating}: a full-run schedule balances sender and
    receiver assignments while preferentially pairing agents who have met less
    often. A recorded pairing seed makes the schedule reproducible.
    \item \textbf{Fixed}: adjacent agents form persistent dyads and roles
    alternate within each dyad. Thus an eight-agent fixed run consists of four
    parallel repeated dyads.
\end{itemize}
Names may be shown or hidden. When names are hidden, prompts use relational
labels such as ``your current co-player'' rather than stable identity labels.
The \texttt{self\_and\_coplayer} history policy can additionally insert a
bounded summary of earlier games involving the current co-player. This is
reputation information assembled from the public simulation record; it is not
access to that co-player's private LLM context.

The completed eight-agent task-order comparison uses balanced rotating
partners, hidden names, and up to three earlier games involving the current
co-player. The completed two-agent comparison is a repeated dyad. These arms
therefore differ simultaneously in population size, partner rotation,
anonymity, and supplied reputation information. They must be described as a
\emph{population-regime comparison}, not as a pure population-size
manipulation. The configured population-isolation follow-up removes these
differences by comparing one fixed dyad with four fixed dyads under unified
prompts, hidden names, and no injected history.

\subsection{Repeated trust game}
\label{sec:protocol-trust-game}

We use a role-swapping repeated trust game \citep{berg1995} with endowment
$E=\$5$ and multiplier $m=3$. In dyad $d$ and round $t$, the sender selects an
actual transfer $s_{td}\in[0,E]$. The receiver's actual allocation is
$m s_{td}$ and the receiver selects an actual return
$r_{td}\in[0,m s_{td}]$. The ground-truth per-round payoffs are
\begin{equation}
    \pi^{S}_{td}=E-s_{td}+r_{td},
    \qquad
    \pi^{R}_{td}=m s_{td}-r_{td}.
    \label{eq:trust-game-payoffs}
\end{equation}
Actual cumulative balances are the sums of these payoffs over the rounds in
which an agent participates.

The sender is prompted first and returns a JSON-like object containing a
\texttt{send} amount. The parsed amount is bounded to $[0,E]$ and stored before
the receiver prompt is constructed. The receiver is then shown the current
communicated send, responds with a \texttt{return} amount, and that amount is
bounded to the feasible actual interval $[0,m s_{td}]$. The complete response,
the parsed action, any clamping, payoffs, and cumulative balances are retained.
A decision response that cannot be parsed causes the run to fail closed rather
than silently inventing an action.

\subsection{Communication noise and the two ledgers}
\label{sec:protocol-noise}

In the primary corrected-v2 condition, both sends and returns are transmitted
through informed bidirectional uniform communication noise with range $a=\$1$.
For an action $x$ with feasible upper bound $u$, the communicated amount is
\begin{equation}
    \widetilde{x}
    = \operatorname{clip}\!\left(x+\epsilon,0,u\right),
    \qquad
    \epsilon\sim\mathcal{U}(-a,a),
    \label{eq:communication-noise}
\end{equation}
rounded to two decimal places. Noise is generated separately for the send and
return events. With seeded protocols, the random draw is a deterministic
function of the noise seed, round, dyad identifier, and action type. This makes
the exogenous communication channel reproducible without coupling it to model
sampling.

The environment maintains parallel actual and communicated views:
\begin{itemize}
    \item The receiver observes $\widetilde{s}_{td}$ and is told that its
    visible receipt is $m\widetilde{s}_{td}$, while the ground-truth allocation
    remains $m s_{td}$.
    \item The sender later observes $\widetilde{r}_{td}$ rather than
    $r_{td}$.
    \item Actual payoffs and balances use Equation~\ref{eq:trust-game-payoffs}.
    The sender's visible payoff is $E-s_{td}+\widetilde{r}_{td}$, and the
    receiver's visible payoff is $m\widetilde{s}_{td}-r_{td}$. Parallel visible
    balances accumulate these quantities for use in later prompts.
\end{itemize}
Thus noise perturbs perceived intentions and feedback, not the underlying
economic transaction. Communicated sends are generated only after the current
actual send exists and while constructing the receiver's prompt. Missing
current transfers raise an error; they never fall back to zero.

\subsection{Myth writing and cultural transmission}
\label{sec:protocol-myths}

In myth-present conditions, every active agent writes one myth per round. In
round one, the prompt requests an approximately 200-word myth and allows any
mythic setting, characters, or symbols. It states that the myth and game are
connected and asks for a story reflecting how the game should be played. In
later rounds, the memory-primary prompt quotes the myth written by the agent's
co-player in the preceding round. It asks the agent to write a new myth that
adapts its earlier work and continues to reflect how the game should be played.
The agent's own earlier myths are available through its private message buffer
rather than duplicated in the current prompt.

Myth transmission therefore has a one-round lag. In a rotating population, an
agent receives the previous round's partner myth, which need not have been
written by its current game partner. The myth channel consequently propagates
text along the realised pairing network rather than creating a population-wide
broadcast. All agents' myth calls within a round are issued concurrently after
their prompts have been prepared, so no agent sees another agent's myth from
the same round.

Before a myth response is committed to memory or the simulation state, a
validator checks that it is non-empty and does not appear to continue into a
game prompt. To avoid false positives from ordinary story language, rejection
requires at least two independent prompt-continuation markers. A rejected myth
is retained in the audit trail, the unanswered prompt is removed from private
memory, and the call is retried once from the same clean conversation state. If
the retry is also invalid, the run stops and writes an error artifact.

\subsection{Scripted-defector and monitor extensions}
\label{sec:protocol-extensions}

The scripted-defector treatment assigns a reproducible subset of agents to a
\texttt{forced\_zero} game policy. A scripted defector always sends zero as a
sender and returns zero as a receiver. These game actions do not consume an LLM
call, but the corresponding prompt and scripted response may still be written
to the agent's private memory so its subsequent myth writing can respond to its
experience. Defectors use the normal LLM myth-writing path. In the planned
primary treatment, two of eight agents are defectors, agent names are hidden,
and the treatment role is not disclosed to the agents.

The engine also contains an optional silent myth monitor. When explicitly
enabled, a separate model flags myths that communicate actionable game
strategy, and a flagged agent's current-round game earnings are removed from
its cumulative balance without a textual notification. The underlying
transaction and payoff records remain unchanged. This monitor is not part of
the corrected-v2 task-order comparison and must not be assumed active unless a
run's metadata records an enabled monitor configuration.

\subsection{Execution, retries, and persistence}
\label{sec:protocol-execution}

At launch, the batch runner records the source commit, whether the worktree is
dirty, the absolute configuration path and SHA-256 digest, provider route,
resolved provider model, maximum output-token setting, task order, prompt arm,
population settings, memory policy, and all explicit seeds. The protected
confirmatory launcher refuses to start from a dirty worktree, uses a fresh
output directory, requires the expected number of completed runs in every
cell, and audits each cell before proceeding.

Provider adapters make up to three attempts for errors they classify as
retryable, with exponential backoff. For the direct Anthropic route used in the
confirmatory batch, this includes rate-limit, connection, timeout, server, and
empty-response failures. If all provider attempts fail, any unanswered user
prompt is removed from the private buffer and the failed attempt remains in the
audit history. Myth validation has the additional single clean retry described
above. Scripted game actions are recorded with a distinct response source so
they are never counted as LLM calls.

After every successful round, the runner atomically writes a lightweight
\texttt{.results.json} containing the simulation state without agent message
histories. At the configured checkpoint interval (ten rounds in the current
batch runners), it also writes a full ordinary checkpoint containing the agent
buffers. If that ordinary checkpoint exists, a subsequent launch resumes at
the first unfinished round. On an exception, the runner writes a separate full
\texttt{.error.json} artifact for diagnosis before re-raising the error; this
error artifact is not selected automatically as a resume point. On successful
completion the runner writes (i) a full-state \texttt{.json}, (ii) the
lightweight results file, (iii) a verbatim \texttt{.log}, and (iv) a
human-readable \texttt{.transcript.pdf}, then removes the ordinary checkpoint.
Output paths are organised by output batch, experiment set, model, task order,
and named game-parameter preset.

\subsection{Protocol audit}
\label{sec:protocol-audit}

Final outputs are admitted to substantive analysis only after a strict audit.
For every run, the audit verifies:
\begin{itemize}
    \item the configured number of rounds, agents, dyads, game responses, myths,
    and accepted interactions;
    \item the absence of unrecovered provider or validation errors and the
    correct accounting of recovered attempts;
    \item the informed-noise notice in system prompts and the numerical bounds
    on every communicated send and return;
    \item the presence of the myth-decision instruction exactly once in each
    applicable current game prompt and its absence from game-only prompts;
    \item private-memory lengths and the configured explicit-history windows;
    \item partner stability in fixed-pair conditions, population-setting text,
    name visibility, scripted-response counts, and defector disclosure policy;
    \item execution provenance, including the source revision, configuration
    digest, provider route, and output-token setting; and
    \item identical realised pairing schedules across runs that share a nominal
    paired-protocol key.
\end{itemize}
The completed corrected-v2 task-order batch contains 60 final runs (six cells,
ten runs per cell). All 60 passed the protocol audit. The completed cells did
not, however, record the intended pairing and noise seeds; their inferential
comparisons must therefore treat runs as independent rather than paired.

\subsection{Outcome measures and statistical unit}
\label{sec:protocol-outcomes}

The independent statistical unit is a complete simulation run, not an
individual dyadic transaction or model call. Primary behavioural summaries are
computed from the actual ledger:
\begin{align}
    \text{trust ratio} &= s_{td}/E, \\
    \text{return ratio} &=
    \begin{cases}
        r_{td}/(m s_{td}), & m s_{td}>0,\\
        0, & m s_{td}=0,
    \end{cases} \\
    \text{final balance} &= \text{mean final actual balance across agents}.
\end{align}
We additionally report mean sent and returned amounts, joint population
balance, and round-level trust and return trajectories. For the completed
six-cell confirmatory comparison, between-condition contrasts use independent-
group Welch tests and Welch confidence intervals, with Holm adjustment across
the six prespecified within-population balance contrasts. Population-regime
interactions are estimated as independent differences-in-differences. Seeded
follow-up designs should preserve replicate identifiers and the shared
exogenous pairing/noise schedules so that their paired structure can be used in
the corresponding analysis.

The linguistic outputs support complementary analyses of within-agent myth
retention, between-agent and population-level semantic convergence, lexical
and structural change, strategy content, and links between myth features and
subsequent game behaviour. Such analyses operate on the accepted myths saved in
the final state; rejected continuation attempts are audit artifacts and are
never treated as cultural outputs.

\subsection{Current experimental matrix and interpretation}
\label{sec:protocol-current-matrix}

Table~\ref{tab:current-protocol-matrix} separates completed evidence from
implemented follow-ups. This distinction is part of the protocol: configured
or smoke-tested cells must not be described as substantive results.

\begin{table}[t]
\centering
\small
\setlength{\tabcolsep}{3pt}
\caption{Current corrected-v2 experimental matrix as of 19 August 2026.}
\label{tab:current-protocol-matrix}
\begin{tabular}{p{0.22\linewidth}p{0.31\linewidth}p{0.15\linewidth}p{0.22\linewidth}}
\toprule
Comparison & Design & Status & Permitted interpretation \\
\midrule
Task order by population regime & Repeated two-agent dyad versus anonymous,
rotating eight-agent population with co-player reputation; game only,
game$\rightarrow$myth, and myth$\rightarrow$game & 60/60 complete and audited
& Task-order and population-regime effects; not a pure agent-count effect \\
\addlinespace
Population isolation & One fixed dyad versus four fixed dyads, unified prompts,
hidden names, and no injected history, crossed with all three task orders &
Implemented; no substantive batch & Cleaner population-size test with partner
structure held fixed \\
\addlinespace
Identity isolation & Eight-agent balanced rotating population, no injected
history, hidden versus stable visible names & Implemented; no substantive batch
& Effect of stable partner identity under rotation \\
\addlinespace
Scripted defectors & Eight-agent rotating population, 0\% versus 25\%
\texttt{forced\_zero} agents, crossed with all three task orders & Implemented
and smoke-tested; no substantive batch & Cultural and behavioural response to
exogenous defection \\
\bottomrule
\end{tabular}
\end{table}

\subsection{Repository sources of truth}
\label{sec:protocol-source-files}

The prose above is implemented by the following repository files:
\begin{itemize}
    \item experiment definitions and prompt text:
    \texttt{config/experiments\_noisy.yaml};
    \item configuration expansion, provenance capture, execution, and output
    layout: \texttt{experiments/run\_noisy\_batch.py};
    \item round orchestration, task ordering, memory modes, checkpoints, and
    state serialization: \texttt{src/simulation.py};
    \item private message buffers and interaction audit records:
    \texttt{src/agents.py};
    \item trust-game prompts, communication noise, actual/visible ledgers, and
    payoff calculations: \texttt{games/trust\_game\_noisy.py};
    \item population assignment, pairings, identity presentation, and scripted
    policies: \texttt{games/dyadic\_pairing.py};
    \item myth prompting, one-round transmission, and response validation:
    \texttt{src/myth\_writer.py};
    \item protected confirmatory launch:
    \texttt{scripts/run\_corrected\_v2\_confirmatory\_20260812.sh}; and
    \item admissibility audit: \texttt{scripts/audit\_v2\_protocol.py}.
\end{itemize}
If this prose and an older manuscript description disagree, the saved run
metadata and these implementation files determine what was actually executed.

% Change \section to \subsection (etc.) if embedding below another heading.
% Requires: amsmath, natbib, booktabs, and array (all already loaded by the
% current Overleaf manuscript). The bibliography should contain berg1995.

\section{Canonical simulation protocol (corrected v2)}
\label{sec:canonical-simulation-protocol}

\paragraph{Scope and precedence.}
This section is the canonical description of the corrected-v2 simulation
workflow. It covers the current memory-primary, informed-noise experiments,
including the completed task-order comparison and the configured population,
identity, and scripted-defector follow-ups. Earlier project phases used
materially different noise and memory implementations and should not be
described using this protocol. In particular, the current implementation keeps
separate actual and communicated ledgers: communication noise changes what an
agent observes about another agent's action, but it does not change the actual
transfer or the ground-truth payoff calculation.

\subsection{System overview}
\label{sec:protocol-overview}

Each simulation run consists of a population of same-model LLM agents playing
ten rounds of a repeated trust game, optionally interleaved with a myth-writing
task. The simulation is centrally orchestrated. Agents do not exchange
free-form chat messages. Their inter-agent information channels are (i) game
actions and communicated outcomes, (ii) any explicitly supplied game-history
or reputation block, and (iii) myths transmitted between rounds. Each agent
has a private rolling chat history; no agent receives another agent's private
conversation memory.

An experiment set is defined in a YAML configuration and expanded into the
Cartesian product of model, system-prompt template, persona, task order, myth
topic, game-parameter preset, myth-prompt arm, and replicate. Each expanded
combination becomes one independent simulation run. Replicates may be executed
in separate worker processes, while the state transitions within any one run
remain ordered. In seeded follow-up protocols, the exogenous pairing and noise
processes are reproducible conditional on their recorded seeds; model sampling
remains stochastic.

The three primary task orders are:
\begin{itemize}
    \item \textbf{Game only}: one game interaction per agent per round and no
    myth-writing task.
    \item \textbf{Game $\rightarrow$ myth}: agents make the current game
    decision before writing the current-round myth. The new myth can therefore
    affect later rounds, but not the game decision that preceded it.
    \item \textbf{Myth $\rightarrow$ game}: agents write the current-round myth
    before making the current game decision. Because both tasks use the same
    private memory, the new myth is available to the immediately following game
    call.
\end{itemize}
Whenever a run contains both tasks, every current game prompt contains exactly
once the instruction that the agent should take myths written in the session
into account when deciding. This explicit link is held constant across the two
myth-present task orders.

\subsection{Agents, prompts, and private memory}
\label{sec:protocol-agents-memory}

All agents in a run share the configured base model, temperature, neutral
persona, and prompt templates, but each agent has a distinct identifier,
optional display name, system prompt, and private message buffer. Unless a
condition states otherwise, current experiments use Claude Sonnet 4.5,
temperature $0.8$, and a neutral persona. The protected confirmatory launcher
routes these calls directly to Anthropic and fixes the maximum response length
at 4,096 tokens.

The system prompt explains both roles of the trust game, the payoff rules, the
required decision format, and the possible myth-writing task. In informed-noise
conditions it also states that communicated amounts may differ slightly from
the amounts actually sent or returned. Population-setting text tells agents
whether partners are fixed or rotating and whether identities are visible.

Memory capacity is measured in completed user--assistant exchanges, not in raw
messages. The system message is always preserved. Before a new model call, the
buffer retains at most the configured number of preceding exchanges and then
adds the current user prompt. In the canonical \emph{memory-primary} regime,
accepted game and myth exchanges are written to the same private buffer. Game-
only cells use capacity three. Two-task cells use capacity six, so both regimes
retain approximately three completed rounds of task interaction: three game
exchanges in the former and roughly three game-plus-myth pairs in the latter.
This matching prevents the addition of the myth task from mechanically
shortening the intended round-level memory horizon.

The exact request and response for every interaction are also copied to a
non-truncated audit history. This audit record includes the prompt, the full
message list supplied to the provider, the response text, available reasoning
and token usage, task and round metadata, role, opponent, dyad, response source,
and any error. The rolling private memory controls agent behaviour; the audit
history is retained only for provenance and verification.

\subsection{Population structure and pairings}
\label{sec:protocol-pairings}

Every round partitions the population into sender--receiver dyads. Within a
dyad, the sender always acts before the receiver. Pairing metadata is saved at
both the round and dyad levels.

For two-agent runs, the same pair remains together for all rounds and sender
and receiver roles alternate by round. For populations larger than two, two
pairing modes are available:
\begin{itemize}
    \item \textbf{Balanced rotating}: a full-run schedule balances sender and
    receiver assignments while preferentially pairing agents who have met less
    often. A recorded pairing seed makes the schedule reproducible.
    \item \textbf{Fixed}: adjacent agents form persistent dyads and roles
    alternate within each dyad. Thus an eight-agent fixed run consists of four
    parallel repeated dyads.
\end{itemize}
Names may be shown or hidden. When names are hidden, prompts use relational
labels such as ``your current co-player'' rather than stable identity labels.
The \texttt{self\_and\_coplayer} history policy can additionally insert a
bounded summary of earlier games involving the current co-player. This is
reputation information assembled from the public simulation record; it is not
access to that co-player's private LLM context.

The completed eight-agent task-order comparison uses balanced rotating
partners, hidden names, and up to three earlier games involving the current
co-player. The completed two-agent comparison is a repeated dyad. These arms
therefore differ simultaneously in population size, partner rotation,
anonymity, and supplied reputation information. They must be described as a
\emph{population-regime comparison}, not as a pure population-size
manipulation. The configured population-isolation follow-up removes these
differences by comparing one fixed dyad with four fixed dyads under unified
prompts, hidden names, and no injected history.

\subsection{Repeated trust game}
\label{sec:protocol-trust-game}

We use a role-swapping repeated trust game \citep{berg1995} with endowment
$E=\$5$ and multiplier $m=3$. In dyad $d$ and round $t$, the sender selects an
actual transfer $s_{td}\in[0,E]$. The receiver's actual allocation is
$m s_{td}$ and the receiver selects an actual return
$r_{td}\in[0,m s_{td}]$. The ground-truth per-round payoffs are
\begin{equation}
    \pi^{S}_{td}=E-s_{td}+r_{td},
    \qquad
    \pi^{R}_{td}=m s_{td}-r_{td}.
    \label{eq:trust-game-payoffs}
\end{equation}
Actual cumulative balances are the sums of these payoffs over the rounds in
which an agent participates.

The sender is prompted first and returns a JSON-like object containing a
\texttt{send} amount. The parsed amount is bounded to $[0,E]$ and stored before
the receiver prompt is constructed. The receiver is then shown the current
communicated send, responds with a \texttt{return} amount, and that amount is
bounded to the feasible actual interval $[0,m s_{td}]$. The complete response,
the parsed action, any clamping, payoffs, and cumulative balances are retained.
A decision response that cannot be parsed causes the run to fail closed rather
than silently inventing an action.

\subsection{Communication noise and the two ledgers}
\label{sec:protocol-noise}

In the primary corrected-v2 condition, both sends and returns are transmitted
through informed bidirectional uniform communication noise with range $a=\$1$.
For an action $x$ with feasible upper bound $u$, the communicated amount is
\begin{equation}
    \widetilde{x}
    = \operatorname{clip}\!\left(x+\epsilon,0,u\right),
    \qquad
    \epsilon\sim\mathcal{U}(-a,a),
    \label{eq:communication-noise}
\end{equation}
rounded to two decimal places. Noise is generated separately for the send and
return events. With seeded protocols, the random draw is a deterministic
function of the noise seed, round, dyad identifier, and action type. This makes
the exogenous communication channel reproducible without coupling it to model
sampling.

The environment maintains parallel actual and communicated views:
\begin{itemize}
    \item The receiver observes $\widetilde{s}_{td}$ and is told that its
    visible receipt is $m\widetilde{s}_{td}$, while the ground-truth allocation
    remains $m s_{td}$.
    \item The sender later observes $\widetilde{r}_{td}$ rather than
    $r_{td}$.
    \item Actual payoffs and balances use Equation~\ref{eq:trust-game-payoffs}.
    The sender's visible payoff is $E-s_{td}+\widetilde{r}_{td}$, and the
    receiver's visible payoff is $m\widetilde{s}_{td}-r_{td}$. Parallel visible
    balances accumulate these quantities for use in later prompts.
\end{itemize}
Thus noise perturbs perceived intentions and feedback, not the underlying
economic transaction. Communicated sends are generated only after the current
actual send exists and while constructing the receiver's prompt. Missing
current transfers raise an error; they never fall back to zero.

\subsection{Myth writing and cultural transmission}
\label{sec:protocol-myths}

In myth-present conditions, every active agent writes one myth per round. In
round one, the prompt requests an approximately 200-word myth and allows any
mythic setting, characters, or symbols. It states that the myth and game are
connected and asks for a story reflecting how the game should be played. In
later rounds, the memory-primary prompt quotes the myth written by the agent's
co-player in the preceding round. It asks the agent to write a new myth that
adapts its earlier work and continues to reflect how the game should be played.
The agent's own earlier myths are available through its private message buffer
rather than duplicated in the current prompt.

Myth transmission therefore has a one-round lag. In a rotating population, an
agent receives the previous round's partner myth, which need not have been
written by its current game partner. The myth channel consequently propagates
text along the realised pairing network rather than creating a population-wide
broadcast. All agents' myth calls within a round are issued concurrently after
their prompts have been prepared, so no agent sees another agent's myth from
the same round.

Before a myth response is committed to memory or the simulation state, a
validator checks that it is non-empty and does not appear to continue into a
game prompt. To avoid false positives from ordinary story language, rejection
requires at least two independent prompt-continuation markers. A rejected myth
is retained in the audit trail, the unanswered prompt is removed from private
memory, and the call is retried once from the same clean conversation state. If
the retry is also invalid, the run stops and writes an error artifact.

\subsection{Scripted-defector and monitor extensions}
\label{sec:protocol-extensions}

The scripted-defector treatment assigns a reproducible subset of agents to a
\texttt{forced\_zero} game policy. A scripted defector always sends zero as a
sender and returns zero as a receiver. These game actions do not consume an LLM
call, but the corresponding prompt and scripted response may still be written
to the agent's private memory so its subsequent myth writing can respond to its
experience. Defectors use the normal LLM myth-writing path. In the planned
primary treatment, two of eight agents are defectors, agent names are hidden,
and the treatment role is not disclosed to the agents.

The engine also contains an optional silent myth monitor. When explicitly
enabled, a separate model flags myths that communicate actionable game
strategy, and a flagged agent's current-round game earnings are removed from
its cumulative balance without a textual notification. The underlying
transaction and payoff records remain unchanged. This monitor is not part of
the corrected-v2 task-order comparison and must not be assumed active unless a
run's metadata records an enabled monitor configuration.

\subsection{Execution, retries, and persistence}
\label{sec:protocol-execution}

At launch, the batch runner records the source commit, whether the worktree is
dirty, the absolute configuration path and SHA-256 digest, provider route,
resolved provider model, maximum output-token setting, task order, prompt arm,
population settings, memory policy, and all explicit seeds. The protected
confirmatory launcher refuses to start from a dirty worktree, uses a fresh
output directory, requires the expected number of completed runs in every
cell, and audits each cell before proceeding.

Provider adapters make up to three attempts for errors they classify as
retryable, with exponential backoff. For the direct Anthropic route used in the
confirmatory batch, this includes rate-limit, connection, timeout, server, and
empty-response failures. If all provider attempts fail, any unanswered user
prompt is removed from the private buffer and the failed attempt remains in the
audit history. Myth validation has the additional single clean retry described
above. Scripted game actions are recorded with a distinct response source so
they are never counted as LLM calls.

After every successful round, the runner atomically writes a lightweight
\texttt{.results.json} containing the simulation state without agent message
histories. At the configured checkpoint interval (ten rounds in the current
batch runners), it also writes a full ordinary checkpoint containing the agent
buffers. If that ordinary checkpoint exists, a subsequent launch resumes at
the first unfinished round. On an exception, the runner writes a separate full
\texttt{.error.json} artifact for diagnosis before re-raising the error; this
error artifact is not selected automatically as a resume point. On successful
completion the runner writes (i) a full-state \texttt{.json}, (ii) the
lightweight results file, (iii) a verbatim \texttt{.log}, and (iv) a
human-readable \texttt{.transcript.pdf}, then removes the ordinary checkpoint.
Output paths are organised by output batch, experiment set, model, task order,
and named game-parameter preset.

\subsection{Protocol audit}
\label{sec:protocol-audit}

Final outputs are admitted to substantive analysis only after a strict audit.
For every run, the audit verifies:
\begin{itemize}
    \item the configured number of rounds, agents, dyads, game responses, myths,
    and accepted interactions;
    \item the absence of unrecovered provider or validation errors and the
    correct accounting of recovered attempts;
    \item the informed-noise notice in system prompts and the numerical bounds
    on every communicated send and return;
    \item the presence of the myth-decision instruction exactly once in each
    applicable current game prompt and its absence from game-only prompts;
    \item private-memory lengths and the configured explicit-history windows;
    \item partner stability in fixed-pair conditions, population-setting text,
    name visibility, scripted-response counts, and defector disclosure policy;
    \item execution provenance, including the source revision, configuration
    digest, provider route, and output-token setting; and
    \item identical realised pairing schedules across runs that share a nominal
    paired-protocol key.
\end{itemize}
The completed corrected-v2 task-order batch contains 60 final runs (six cells,
ten runs per cell). All 60 passed the protocol audit. The completed cells did
not, however, record the intended pairing and noise seeds; their inferential
comparisons must therefore treat runs as independent rather than paired.

\subsection{Outcome measures and statistical unit}
\label{sec:protocol-outcomes}

The independent statistical unit is a complete simulation run, not an
individual dyadic transaction or model call. Primary behavioural summaries are
computed from the actual ledger:
\begin{align}
    \text{trust ratio} &= s_{td}/E, \\
    \text{return ratio} &=
    \begin{cases}
        r_{td}/(m s_{td}), & m s_{td}>0,\\
        0, & m s_{td}=0,
    \end{cases} \\
    \text{final balance} &= \text{mean final actual balance across agents}.
\end{align}
We additionally report mean sent and returned amounts, joint population
balance, and round-level trust and return trajectories. For the completed
six-cell confirmatory comparison, between-condition contrasts use independent-
group Welch tests and Welch confidence intervals, with Holm adjustment across
the six prespecified within-population balance contrasts. Population-regime
interactions are estimated as independent differences-in-differences. Seeded
follow-up designs should preserve replicate identifiers and the shared
exogenous pairing/noise schedules so that their paired structure can be used in
the corresponding analysis.

The linguistic outputs support complementary analyses of within-agent myth
retention, between-agent and population-level semantic convergence, lexical
and structural change, strategy content, and links between myth features and
subsequent game behaviour. Such analyses operate on the accepted myths saved in
the final state; rejected continuation attempts are audit artifacts and are
never treated as cultural outputs.

\subsection{Current experimental matrix and interpretation}
\label{sec:protocol-current-matrix}

Table~\ref{tab:current-protocol-matrix} separates completed evidence from
implemented follow-ups. This distinction is part of the protocol: configured
or smoke-tested cells must not be described as substantive results.

\begin{table}[t]
\centering
\small
\setlength{\tabcolsep}{3pt}
\caption{Current corrected-v2 experimental matrix as of 19 August 2026.}
\label{tab:current-protocol-matrix}
\begin{tabular}{p{0.22\linewidth}p{0.31\linewidth}p{0.15\linewidth}p{0.22\linewidth}}
\toprule
Comparison & Design & Status & Permitted interpretation \\
\midrule
Task order by population regime & Repeated two-agent dyad versus anonymous,
rotating eight-agent population with co-player reputation; game only,
game$\rightarrow$myth, and myth$\rightarrow$game & 60/60 complete and audited
& Task-order and population-regime effects; not a pure agent-count effect \\
\addlinespace
Population isolation & One fixed dyad versus four fixed dyads, unified prompts,
hidden names, and no injected history, crossed with all three task orders &
Implemented; no substantive batch & Cleaner population-size test with partner
structure held fixed \\
\addlinespace
Identity isolation & Eight-agent balanced rotating population, no injected
history, hidden versus stable visible names & Implemented; no substantive batch
& Effect of stable partner identity under rotation \\
\addlinespace
Scripted defectors & Eight-agent rotating population, 0\% versus 25\%
\texttt{forced\_zero} agents, crossed with all three task orders & Implemented
and smoke-tested; no substantive batch & Cultural and behavioural response to
exogenous defection \\
\bottomrule
\end{tabular}
\end{table}

\subsection{Repository sources of truth}
\label{sec:protocol-source-files}

The prose above is implemented by the following repository files:
\begin{itemize}
    \item experiment definitions and prompt text:
    \texttt{config/experiments\_noisy.yaml};
    \item configuration expansion, provenance capture, execution, and output
    layout: \texttt{experiments/run\_noisy\_batch.py};
    \item round orchestration, task ordering, memory modes, checkpoints, and
    state serialization: \texttt{src/simulation.py};
    \item private message buffers and interaction audit records:
    \texttt{src/agents.py};
    \item trust-game prompts, communication noise, actual/visible ledgers, and
    payoff calculations: \texttt{games/trust\_game\_noisy.py};
    \item population assignment, pairings, identity presentation, and scripted
    policies: \texttt{games/dyadic\_pairing.py};
    \item myth prompting, one-round transmission, and response validation:
    \texttt{src/myth\_writer.py};
    \item protected confirmatory launch:
    \texttt{scripts/run\_corrected\_v2\_confirmatory\_20260812.sh}; and
    \item admissibility audit: \texttt{scripts/audit\_v2\_protocol.py}.
\end{itemize}
If this prose and an older manuscript description disagree, the saved run
metadata and these implementation files determine what was actually executed.

% Change \section to \subsection (etc.) if embedding below another heading.
% Requires: amsmath, natbib, booktabs, and array (all already loaded by the
% current Overleaf manuscript). The bibliography should contain berg1995.

\section{Canonical simulation protocol (corrected v2)}
\label{sec:canonical-simulation-protocol}

\paragraph{Scope and precedence.}
This section is the canonical description of the corrected-v2 simulation
workflow. It covers the current memory-primary, informed-noise experiments,
including the completed task-order comparison and the configured population,
identity, and scripted-defector follow-ups. Earlier project phases used
materially different noise and memory implementations and should not be
described using this protocol. In particular, the current implementation keeps
separate actual and communicated ledgers: communication noise changes what an
agent observes about another agent's action, but it does not change the actual
transfer or the ground-truth payoff calculation.

\subsection{System overview}
\label{sec:protocol-overview}

Each simulation run consists of a population of same-model LLM agents playing
ten rounds of a repeated trust game, optionally interleaved with a myth-writing
task. The simulation is centrally orchestrated. Agents do not exchange
free-form chat messages. Their inter-agent information channels are (i) game
actions and communicated outcomes, (ii) any explicitly supplied game-history
or reputation block, and (iii) myths transmitted between rounds. Each agent
has a private rolling chat history; no agent receives another agent's private
conversation memory.

An experiment set is defined in a YAML configuration and expanded into the
Cartesian product of model, system-prompt template, persona, task order, myth
topic, game-parameter preset, myth-prompt arm, and replicate. Each expanded
combination becomes one independent simulation run. Replicates may be executed
in separate worker processes, while the state transitions within any one run
remain ordered. In seeded follow-up protocols, the exogenous pairing and noise
processes are reproducible conditional on their recorded seeds; model sampling
remains stochastic.

The three primary task orders are:
\begin{itemize}
    \item \textbf{Game only}: one game interaction per agent per round and no
    myth-writing task.
    \item \textbf{Game $\rightarrow$ myth}: agents make the current game
    decision before writing the current-round myth. The new myth can therefore
    affect later rounds, but not the game decision that preceded it.
    \item \textbf{Myth $\rightarrow$ game}: agents write the current-round myth
    before making the current game decision. Because both tasks use the same
    private memory, the new myth is available to the immediately following game
    call.
\end{itemize}
Whenever a run contains both tasks, every current game prompt contains exactly
once the instruction that the agent should take myths written in the session
into account when deciding. This explicit link is held constant across the two
myth-present task orders.

\subsection{Agents, prompts, and private memory}
\label{sec:protocol-agents-memory}

All agents in a run share the configured base model, temperature, neutral
persona, and prompt templates, but each agent has a distinct identifier,
optional display name, system prompt, and private message buffer. Unless a
condition states otherwise, current experiments use Claude Sonnet 4.5,
temperature $0.8$, and a neutral persona. The protected confirmatory launcher
routes these calls directly to Anthropic and fixes the maximum response length
at 4,096 tokens.

The system prompt explains both roles of the trust game, the payoff rules, the
required decision format, and the possible myth-writing task. In informed-noise
conditions it also states that communicated amounts may differ slightly from
the amounts actually sent or returned. Population-setting text tells agents
whether partners are fixed or rotating and whether identities are visible.

Memory capacity is measured in completed user--assistant exchanges, not in raw
messages. The system message is always preserved. Before a new model call, the
buffer retains at most the configured number of preceding exchanges and then
adds the current user prompt. In the canonical \emph{memory-primary} regime,
accepted game and myth exchanges are written to the same private buffer. Game-
only cells use capacity three. Two-task cells use capacity six, so both regimes
retain approximately three completed rounds of task interaction: three game
exchanges in the former and roughly three game-plus-myth pairs in the latter.
This matching prevents the addition of the myth task from mechanically
shortening the intended round-level memory horizon.

The exact request and response for every interaction are also copied to a
non-truncated audit history. This audit record includes the prompt, the full
message list supplied to the provider, the response text, available reasoning
and token usage, task and round metadata, role, opponent, dyad, response source,
and any error. The rolling private memory controls agent behaviour; the audit
history is retained only for provenance and verification.

\subsection{Population structure and pairings}
\label{sec:protocol-pairings}

Every round partitions the population into sender--receiver dyads. Within a
dyad, the sender always acts before the receiver. Pairing metadata is saved at
both the round and dyad levels.

For two-agent runs, the same pair remains together for all rounds and sender
and receiver roles alternate by round. For populations larger than two, two
pairing modes are available:
\begin{itemize}
    \item \textbf{Balanced rotating}: a full-run schedule balances sender and
    receiver assignments while preferentially pairing agents who have met less
    often. A recorded pairing seed makes the schedule reproducible.
    \item \textbf{Fixed}: adjacent agents form persistent dyads and roles
    alternate within each dyad. Thus an eight-agent fixed run consists of four
    parallel repeated dyads.
\end{itemize}
Names may be shown or hidden. When names are hidden, prompts use relational
labels such as ``your current co-player'' rather than stable identity labels.
The \texttt{self\_and\_coplayer} history policy can additionally insert a
bounded summary of earlier games involving the current co-player. This is
reputation information assembled from the public simulation record; it is not
access to that co-player's private LLM context.

The completed eight-agent task-order comparison uses balanced rotating
partners, hidden names, and up to three earlier games involving the current
co-player. The completed two-agent comparison is a repeated dyad. These arms
therefore differ simultaneously in population size, partner rotation,
anonymity, and supplied reputation information. They must be described as a
\emph{population-regime comparison}, not as a pure population-size
manipulation. The configured population-isolation follow-up removes these
differences by comparing one fixed dyad with four fixed dyads under unified
prompts, hidden names, and no injected history.

\subsection{Repeated trust game}
\label{sec:protocol-trust-game}

We use a role-swapping repeated trust game \citep{berg1995} with endowment
$E=\$5$ and multiplier $m=3$. In dyad $d$ and round $t$, the sender selects an
actual transfer $s_{td}\in[0,E]$. The receiver's actual allocation is
$m s_{td}$ and the receiver selects an actual return
$r_{td}\in[0,m s_{td}]$. The ground-truth per-round payoffs are
\begin{equation}
    \pi^{S}_{td}=E-s_{td}+r_{td},
    \qquad
    \pi^{R}_{td}=m s_{td}-r_{td}.
    \label{eq:trust-game-payoffs}
\end{equation}
Actual cumulative balances are the sums of these payoffs over the rounds in
which an agent participates.

The sender is prompted first and returns a JSON-like object containing a
\texttt{send} amount. The parsed amount is bounded to $[0,E]$ and stored before
the receiver prompt is constructed. The receiver is then shown the current
communicated send, responds with a \texttt{return} amount, and that amount is
bounded to the feasible actual interval $[0,m s_{td}]$. The complete response,
the parsed action, any clamping, payoffs, and cumulative balances are retained.
A decision response that cannot be parsed causes the run to fail closed rather
than silently inventing an action.

\subsection{Communication noise and the two ledgers}
\label{sec:protocol-noise}

In the primary corrected-v2 condition, both sends and returns are transmitted
through informed bidirectional uniform communication noise with range $a=\$1$.
For an action $x$ with feasible upper bound $u$, the communicated amount is
\begin{equation}
    \widetilde{x}
    = \operatorname{clip}\!\left(x+\epsilon,0,u\right),
    \qquad
    \epsilon\sim\mathcal{U}(-a,a),
    \label{eq:communication-noise}
\end{equation}
rounded to two decimal places. Noise is generated separately for the send and
return events. With seeded protocols, the random draw is a deterministic
function of the noise seed, round, dyad identifier, and action type. This makes
the exogenous communication channel reproducible without coupling it to model
sampling.

The environment maintains parallel actual and communicated views:
\begin{itemize}
    \item The receiver observes $\widetilde{s}_{td}$ and is told that its
    visible receipt is $m\widetilde{s}_{td}$, while the ground-truth allocation
    remains $m s_{td}$.
    \item The sender later observes $\widetilde{r}_{td}$ rather than
    $r_{td}$.
    \item Actual payoffs and balances use Equation~\ref{eq:trust-game-payoffs}.
    The sender's visible payoff is $E-s_{td}+\widetilde{r}_{td}$, and the
    receiver's visible payoff is $m\widetilde{s}_{td}-r_{td}$. Parallel visible
    balances accumulate these quantities for use in later prompts.
\end{itemize}
Thus noise perturbs perceived intentions and feedback, not the underlying
economic transaction. Communicated sends are generated only after the current
actual send exists and while constructing the receiver's prompt. Missing
current transfers raise an error; they never fall back to zero.

\subsection{Myth writing and cultural transmission}
\label{sec:protocol-myths}

In myth-present conditions, every active agent writes one myth per round. In
round one, the prompt requests an approximately 200-word myth and allows any
mythic setting, characters, or symbols. It states that the myth and game are
connected and asks for a story reflecting how the game should be played. In
later rounds, the memory-primary prompt quotes the myth written by the agent's
co-player in the preceding round. It asks the agent to write a new myth that
adapts its earlier work and continues to reflect how the game should be played.
The agent's own earlier myths are available through its private message buffer
rather than duplicated in the current prompt.

Myth transmission therefore has a one-round lag. In a rotating population, an
agent receives the previous round's partner myth, which need not have been
written by its current game partner. The myth channel consequently propagates
text along the realised pairing network rather than creating a population-wide
broadcast. All agents' myth calls within a round are issued concurrently after
their prompts have been prepared, so no agent sees another agent's myth from
the same round.

Before a myth response is committed to memory or the simulation state, a
validator checks that it is non-empty and does not appear to continue into a
game prompt. To avoid false positives from ordinary story language, rejection
requires at least two independent prompt-continuation markers. A rejected myth
is retained in the audit trail, the unanswered prompt is removed from private
memory, and the call is retried once from the same clean conversation state. If
the retry is also invalid, the run stops and writes an error artifact.

\subsection{Scripted-defector and monitor extensions}
\label{sec:protocol-extensions}

The scripted-defector treatment assigns a reproducible subset of agents to a
\texttt{forced\_zero} game policy. A scripted defector always sends zero as a
sender and returns zero as a receiver. These game actions do not consume an LLM
call, but the corresponding prompt and scripted response may still be written
to the agent's private memory so its subsequent myth writing can respond to its
experience. Defectors use the normal LLM myth-writing path. In the planned
primary treatment, two of eight agents are defectors, agent names are hidden,
and the treatment role is not disclosed to the agents.

The engine also contains an optional silent myth monitor. When explicitly
enabled, a separate model flags myths that communicate actionable game
strategy, and a flagged agent's current-round game earnings are removed from
its cumulative balance without a textual notification. The underlying
transaction and payoff records remain unchanged. This monitor is not part of
the corrected-v2 task-order comparison and must not be assumed active unless a
run's metadata records an enabled monitor configuration.

\subsection{Execution, retries, and persistence}
\label{sec:protocol-execution}

At launch, the batch runner records the source commit, whether the worktree is
dirty, the absolute configuration path and SHA-256 digest, provider route,
resolved provider model, maximum output-token setting, task order, prompt arm,
population settings, memory policy, and all explicit seeds. The protected
confirmatory launcher refuses to start from a dirty worktree, uses a fresh
output directory, requires the expected number of completed runs in every
cell, and audits each cell before proceeding.

Provider adapters make up to three attempts for errors they classify as
retryable, with exponential backoff. For the direct Anthropic route used in the
confirmatory batch, this includes rate-limit, connection, timeout, server, and
empty-response failures. If all provider attempts fail, any unanswered user
prompt is removed from the private buffer and the failed attempt remains in the
audit history. Myth validation has the additional single clean retry described
above. Scripted game actions are recorded with a distinct response source so
they are never counted as LLM calls.

After every successful round, the runner atomically writes a lightweight
\texttt{.results.json} containing the simulation state without agent message
histories. At the configured checkpoint interval (ten rounds in the current
batch runners), it also writes a full ordinary checkpoint containing the agent
buffers. If that ordinary checkpoint exists, a subsequent launch resumes at
the first unfinished round. On an exception, the runner writes a separate full
\texttt{.error.json} artifact for diagnosis before re-raising the error; this
error artifact is not selected automatically as a resume point. On successful
completion the runner writes (i) a full-state \texttt{.json}, (ii) the
lightweight results file, (iii) a verbatim \texttt{.log}, and (iv) a
human-readable \texttt{.transcript.pdf}, then removes the ordinary checkpoint.
Output paths are organised by output batch, experiment set, model, task order,
and named game-parameter preset.

\subsection{Protocol audit}
\label{sec:protocol-audit}

Final outputs are admitted to substantive analysis only after a strict audit.
For every run, the audit verifies:
\begin{itemize}
    \item the configured number of rounds, agents, dyads, game responses, myths,
    and accepted interactions;
    \item the absence of unrecovered provider or validation errors and the
    correct accounting of recovered attempts;
    \item the informed-noise notice in system prompts and the numerical bounds
    on every communicated send and return;
    \item the presence of the myth-decision instruction exactly once in each
    applicable current game prompt and its absence from game-only prompts;
    \item private-memory lengths and the configured explicit-history windows;
    \item partner stability in fixed-pair conditions, population-setting text,
    name visibility, scripted-response counts, and defector disclosure policy;
    \item execution provenance, including the source revision, configuration
    digest, provider route, and output-token setting; and
    \item identical realised pairing schedules across runs that share a nominal
    paired-protocol key.
\end{itemize}
The completed corrected-v2 task-order batch contains 60 final runs (six cells,
ten runs per cell). All 60 passed the protocol audit. The completed cells did
not, however, record the intended pairing and noise seeds; their inferential
comparisons must therefore treat runs as independent rather than paired.

\subsection{Outcome measures and statistical unit}
\label{sec:protocol-outcomes}

The independent statistical unit is a complete simulation run, not an
individual dyadic transaction or model call. Primary behavioural summaries are
computed from the actual ledger:
\begin{align}
    \text{trust ratio} &= s_{td}/E, \\
    \text{return ratio} &=
    \begin{cases}
        r_{td}/(m s_{td}), & m s_{td}>0,\\
        0, & m s_{td}=0,
    \end{cases} \\
    \text{final balance} &= \text{mean final actual balance across agents}.
\end{align}
We additionally report mean sent and returned amounts, joint population
balance, and round-level trust and return trajectories. For the completed
six-cell confirmatory comparison, between-condition contrasts use independent-
group Welch tests and Welch confidence intervals, with Holm adjustment across
the six prespecified within-population balance contrasts. Population-regime
interactions are estimated as independent differences-in-differences. Seeded
follow-up designs should preserve replicate identifiers and the shared
exogenous pairing/noise schedules so that their paired structure can be used in
the corresponding analysis.

The linguistic outputs support complementary analyses of within-agent myth
retention, between-agent and population-level semantic convergence, lexical
and structural change, strategy content, and links between myth features and
subsequent game behaviour. Such analyses operate on the accepted myths saved in
the final state; rejected continuation attempts are audit artifacts and are
never treated as cultural outputs.

\subsection{Current experimental matrix and interpretation}
\label{sec:protocol-current-matrix}

Table~\ref{tab:current-protocol-matrix} separates completed evidence from
implemented follow-ups. This distinction is part of the protocol: configured
or smoke-tested cells must not be described as substantive results.

\begin{table}[t]
\centering
\small
\setlength{\tabcolsep}{3pt}
\caption{Current corrected-v2 experimental matrix as of 19 August 2026.}
\label{tab:current-protocol-matrix}
\begin{tabular}{p{0.22\linewidth}p{0.31\linewidth}p{0.15\linewidth}p{0.22\linewidth}}
\toprule
Comparison & Design & Status & Permitted interpretation \\
\midrule
Task order by population regime & Repeated two-agent dyad versus anonymous,
rotating eight-agent population with co-player reputation; game only,
game$\rightarrow$myth, and myth$\rightarrow$game & 60/60 complete and audited
& Task-order and population-regime effects; not a pure agent-count effect \\
\addlinespace
Population isolation & One fixed dyad versus four fixed dyads, unified prompts,
hidden names, and no injected history, crossed with all three task orders &
Implemented; no substantive batch & Cleaner population-size test with partner
structure held fixed \\
\addlinespace
Identity isolation & Eight-agent balanced rotating population, no injected
history, hidden versus stable visible names & Implemented; no substantive batch
& Effect of stable partner identity under rotation \\
\addlinespace
Scripted defectors & Eight-agent rotating population, 0\% versus 25\%
\texttt{forced\_zero} agents, crossed with all three task orders & Implemented
and smoke-tested; no substantive batch & Cultural and behavioural response to
exogenous defection \\
\bottomrule
\end{tabular}
\end{table}

\subsection{Repository sources of truth}
\label{sec:protocol-source-files}

The prose above is implemented by the following repository files:
\begin{itemize}
    \item experiment definitions and prompt text:
    \texttt{config/experiments\_noisy.yaml};
    \item configuration expansion, provenance capture, execution, and output
    layout: \texttt{experiments/run\_noisy\_batch.py};
    \item round orchestration, task ordering, memory modes, checkpoints, and
    state serialization: \texttt{src/simulation.py};
    \item private message buffers and interaction audit records:
    \texttt{src/agents.py};
    \item trust-game prompts, communication noise, actual/visible ledgers, and
    payoff calculations: \texttt{games/trust\_game\_noisy.py};
    \item population assignment, pairings, identity presentation, and scripted
    policies: \texttt{games/dyadic\_pairing.py};
    \item myth prompting, one-round transmission, and response validation:
    \texttt{src/myth\_writer.py};
    \item protected confirmatory launch:
    \texttt{scripts/run\_corrected\_v2\_confirmatory\_20260812.sh}; and
    \item admissibility audit: \texttt{scripts/audit\_v2\_protocol.py}.
\end{itemize}
If this prose and an older manuscript description disagree, the saved run
metadata and these implementation files determine what was actually executed.
